\section{Conclusion}
\label{sec:conclusion}

We presented a parallel five-track corpus of 13{,}077 Sri Lankan banking support
tickets and used its structure to locate where multilingual pretraining stops
transferring. Because the same tickets appear in native script and in
transliteration, the comparison isolates orthography from language, and the
encoder's advantage over a character $n$-gram baseline lives entirely on scripts
seen during pretraining: $+0.0698$ as a difference of differences, with the same
shape reproduced between two encoders and on a second task. The mechanism is
tokenizer vocabulary coverage rather than sequence fragmentation, which makes it
diagnosable before training and unfixable by fine-tuning.

Applying published scheduling rules to the resulting posteriors, the ordering
rule moves queueing cost by a factor of three where the probability model moves
it by 3\%, and modelling dependence among the three labels improves calibration
by 27\% without improving accuracy. A study reporting only \macrof{} would have
missed both.

The corpus, the frozen split, and the model checkpoints are publicly available
under CC BY 4.0 with attribution to BANKING77.
\ifblind\relax\note{Repository URLs withheld for double-blind review; add for
camera-ready.}\fi
